\section{Separating teacher precision from extra learner depth}
\label{sec:bounded-reset-padding}
For ReLU and a clipped observation, the scalar output can be clipped
inside the comparison network before the extra layers. The two
carried features are then bounded, even if the teacher coefficients
are large. Different relay layers can use different comparison
precisions. Their cost need not multiply the teacher coefficient
logarithm by learner depth.

\begin{theorem}[Bounded reset and layer-specific relay precision]
\label{thm:bounded-reset-padding}
Let a ReLU teacher have hidden depth $h$, width $m$, coefficient cap
$B\ge1$, and rooted formula size $S$. Its observed output is
$g=s(z_g)$, where $s$ clips to $[-1,1]$.
Use the fixed small-intensity L\'evy fan-in prior, and prescribe any
learner depth $H\ge h$ and width $N\ge\max(m+1,3)$.
Put
\[
 L=(h+3)\log\{A(d+1)(m+1)(B+8)\},
 \qquad S_0=S+h+1,
\]
and
\begin{align}
 \mathcal J_n^{\mathrm{reset}}
 =A\bigl\{&(d+1)S_0
       [L+\log\{A(S_0+H+1)(n+1)\}]\notag\\
 &+H\log\{A(S_0+H+1)(H+1)(n+1)\}+L^2\bigr\}.
 \label{eq:reset-resource}
\end{align}
For the same conditionally sub-Gaussian label assumptions and fixed
likelihood temperature as the sharp all-data expectation theorem,
\begin{equation}
 \E R(\bar f_{S_n})
 \le R(g)+\frac{A(1+\sigma^2)(1+\mathcal J_n^{\mathrm{reset}})}{n+1}.
 \label{eq:reset-oracle}
\end{equation}
The constant $A$ depends only on the fixed L\'evy law. In particular,
the term proportional to extra learner depth contains no factor
$\log B$. Large teacher coefficients enter through the teacher term
and the additive $L^2$ term.
\end{theorem}
\begin{proof}
We give a data-independent comparison law with expected uniform raw
error at most $r$ and entropy bounded by the right side of
\eqref{eq:reset-resource}, replacing $n+1$ by $1/r$.
The expected-risk result then follows with $r=(n+1)^{-1}$.

\emph{An exact bounded reset.}
Carry a constant-one feature through the original $h$ layers, adding
at most one neuron per layer. After the teacher computation add
\[
 a=\operatorname{ReLU}(z_g+1),\quad
 b=\operatorname{ReLU}(z_g-1),\quad u=1,
\]
and then
\[
 q_+=\operatorname{ReLU}(a-b-u),\qquad
 q_-=\operatorname{ReLU}(-a+b+u).
\]
Because $a-b-u=s(z_g)$, these two features lie in $[0,1]$ and
$q_+-q_-=g$. This uses two extra hidden layers, width at most
$\max(m+1,3)$, coefficients bounded by $B+1$, and rooted formula
size at most $A(S+h+1)$. The constant paths and the fixed number of
copies of the teacher formula account for this count.
If $H<h+2$, use the original comparison directly; since the number
of extra layers is at most one, its entropy is already bounded by
\eqref{eq:reset-resource}. Thus assume $H\ge h+2$ and write
$q=H-h-2$.

\emph{Teacher block.}
Apply the existing uniform comparison construction through the reset
block at a common precision $\eta_0=r e^{-A L}$, with a sufficiently
large numerical multiplier in $A L$. Its expected uniform error at
each selected reset feature is at most $r/100$. Its selected-feature
first moments are therefore at most two. Its average native-and-selected
first moment is at most some deterministic $U\ge1$ with
\begin{equation}
 \log U\le A L.
 \label{eq:reset-bulk-envelope}
\end{equation}
These are the same absolute-path first-moment bounds used in the
uniform comparison proof. The entropy paid by the teacher block is
at most
\[
 A(d+1)(S+h+1)
       [L+\log\{A(S+h+q+1)/r\}].
\]
The group-size logarithm already permits all later relay occurrences.

\emph{Relay block.}
Let $T_\ell$ be the largest expected uniform absolute value of either
selected relay feature, and let
$P_\ell=N^{-1}\sum_j\E\sup_x|h_j^\ell(x)|$.
At the end of the reset, $T_0\le2$ and $P_0\le U$.
Write $\rho=\int|u|\mathfrak n(du)\le1/4$.
For $q\ge1$, choose deterministic layer precisions
\begin{equation}
 \eta_\ell=\frac{c r}{q[1+U\rho^{\ell-1}]},
 \qquad 1\le\ell\le q,
 \label{eq:variable-relay-precision}
\end{equation}
with a sufficiently small fixed $c$. They are at most $1/8$.
The designated coefficient-one law then has positive support and
mean one, hence expected absolute value one. Its expected deviation
from one is at most $A\eta_\ell$. Background coordinates have
expected absolute value at most $A\eta_\ell/N$. Unselected rows
remain native. Independence from preceding features gives
\begin{align}
 T_\ell&\le T_{\ell-1}+A\eta_\ell P_{\ell-1},\notag\\
 P_\ell&\le\rho P_{\ell-1}+2T_\ell.
 \label{eq:variable-relay-moments}
\end{align}
Induction, using \eqref{eq:variable-relay-precision}, gives
\[
 T_\ell\le3,\qquad P_\ell\le8+U\rho^\ell.
\]
Indeed, the second bound follows from the first and the geometric
sum, and the total increase in the first is at most
$A\sum_\ell\eta_\ell(8+U\rho^{\ell-1})\le8Acr\le1$.

If $e_\ell$ is the maximum expected uniform error from the ideal
reset feature, Lipschitzness and independence give
\[
 e_\ell\le e_{\ell-1}
       +A\eta_\ell(T_{\ell-1}+P_{\ell-1}).
\]
The total error increase is at most $Acr$. Use the final readout
coefficients $1,-1$ at precision
$c r/[1+U\rho^q]$ and the corresponding background law.
The readout adds another $O(cr)$ error. After choosing $c$ small,
the expected uniform raw error is at most $r$. For $q=0$, apply the
same readout prescription directly after the reset.

\emph{Entropy and width cancellation.}
The proof of the disjoint-embedding entropy identity works with a
different known cutoff at each layer: the parent row's cutoff
identifies each designated incoming edge. Thus the $\log N$ costs
still cancel against embedding location entropy. In the small-width
branch, the shared reset-and-relay network fits, and $\log N$ is
bounded by the logarithm of its total formula resource.
There are two selected relay rows per layer, each with one designated
coefficient bounded by one. Their total entropy, including the final
readout, is at most
\begin{equation}
 A(q+1)\log\frac{A(S+h+q+1)(q+1)}r
       +A\sum_{\ell=0}^{q}\log(1+U\rho^\ell).
 \label{eq:relay-entropy-sum}
\end{equation}
The elementary geometric estimate
\begin{equation}
 \sum_{\ell\ge0}\log(1+U\rho^\ell)
       \le A_\rho[1+\log^2(1+U)]
 \label{eq:geometric-log-sum}
\end{equation}
follows by splitting at
$\ell=\lceil\log U/\log(1/\rho)\rceil$: before this index use
$\log(1+v)\le\log2+\log v$ for $v\ge1$, and afterward use
$\log(1+v)\le v$. Combine
\eqref{eq:reset-bulk-envelope}--\eqref{eq:geometric-log-sum}
with the teacher-block entropy. This is the resource
\eqref{eq:reset-resource} with $n+1=1/r$.
The observation clipping cannot enlarge the uniform error. Since
the conditional mean and observed predictions lie in $[-1,1]$,
comparison risk increases by at most $4r$. The sharp all-data
expectation theorem proves \eqref{eq:reset-oracle}.
\end{proof}

\paragraph{Why a fixed tiny cutoff is unnecessary.}
The large teacher block may leave a large average native background
$U$. It contracts geometrically in subsequent untouched rows. Only
the early relay layers need precision inversely proportional to this
background. The sum of those precision logarithms is quadratic in
$\log U$, independent of the number of later layers. The prior itself
is unchanged; all layer-specific cutoffs belong to a comparison law.


\paragraph{Width conventions.}
If the teacher already carries a constant-one feature through its
hidden layers, the extra width-one allowance is unnecessary: the
construction uses width at most $\max(m,3)$. If its raw output is
already bounded by one, the two-channel compiler itself is a bounded
reset; no constant wire or two-stage clipping gadget is needed, and
width $N\ge m$ suffices (using the width-one special case when needed).
These variants have the same resource bound, with adjusted numerical
constants.
